% 题干转写：完成
% 解析转写：完成
% 图形核对：完成
\section{第38届复赛第5题：原子激光制冷与塞曼减速}

\subsection{题干}
原子激光制冷利用与原子束反向传播的单频激光，使原子反复吸收光子并各向同性发射，从而减小原子束流速度。塞曼减速装置利用空间变化的磁场补偿原子减速时的多普勒失谐，使原子在减速管中保持共振散射，如图\ref{fig:cpho38q05}所示。

考虑经准直的 ${}^{40}\mathrm K$ 原子束，其初始运动只有一个自由度，原子气体温度为 $350\,{}^\circ\mathrm C$。共振时散射速率 $\Gamma=5.00\times10^6\,\mathrm{s^{-1}}$，减速激光波长 $\lambda=670\,\mathrm{nm}$。跃迁频率随磁感应强度满足
\[
f_0(B)=f_0(B=0)+\beta B,\qquad
\beta=-1.00\times10^{10}\,\mathrm{Hz/T}.
\]
求：（1）原子的方均根速率；（2）共振减速的加速度和减速至静止的距离；（3）无磁场时使相向运动原子共振所需的激光降频量；（4）以准直管出口 $z=0$ 处 $B=0$ 为条件，实现沿程共振的 $B(z)$。
\begin{figure}[htbp]
  \centering
  \begin{tikzpicture}[x=.85cm,y=.85cm,>=Stealth,font=\small]
  % Atomic oven and collimation aperture.
  \draw[thick,fill=gray!20] (0,0) rectangle (1.45,1.35);
  \draw[thick,fill=gray!55] (0,0) rectangle (1.45,.23);
  \foreach \x/\y in {.25/.52,.62/.78,1.04/.49,.36/1.06,.92/1.08,1.21/.83}
    \fill (\x,\y) circle (1.3pt);
  \node[below] at (.72,0) {原子炉};
  \draw[thick] (1.45,.76)--(2.12,.76);
  \draw[thick] (1.45,.88)--(2.12,.88);
  \draw[thick] (2.12,.56)--(2.12,1.08);
  \draw[thick] (2.28,.63)--(2.28,1.01);
  \node[above] at (2.2,1.25) {准直管};

  % Atomic beam and deceleration tube.
  \draw[thick] (2.28,.63)--(3.1,.63)--(3.45,.42)--(6.7,.42);
  \draw[thick] (2.28,1.01)--(3.1,1.01)--(3.45,1.22)--(6.7,1.22);
  \draw[->,very thick] (2.55,.82)--(4.25,.82) node[midway,above] {原子束流};
  \node[above] at (4.95,1.25) {减速管};

  % Solenoid around the deceleration tube.
  \foreach \x in {3.55,3.82,...,6.45}
    \draw[gray!70,thick] (\x,.22) -- (\x,1.42);
  \draw[gray!70,thick] (3.38,.22)--(6.68,.22);
  \draw[gray!70,thick] (3.38,1.42)--(6.68,1.42);
  \node[below] at (5.0,.05) {线圈（提供均匀磁场）};

  % Cold-atom experimental chamber and counter-propagating laser.
  \draw[thick] (6.7,.25) rectangle (8.25,1.39);
  \node[above] at (7.48,1.39) {冷原子实验腔};
  \draw[thick] (8.25,.67)--(9.25,.67);
  \draw[thick] (8.25,.97)--(9.25,.97);
  \draw[->,very thick] (9.42,.82)--(6.93,.82)
    node[midway,below] {减速激光};
\end{tikzpicture}

  \caption{第38届复赛第5题塞曼减速装置的TikZ重绘。}
  \label{fig:cpho38q05}
\end{figure}

\subsection{参考解答整理}
\paragraph{（1）准直原子的初始速率。}
准直后只保留一个自由度，能量均分给出
\[
\frac12mv_0^2=\frac12k_\mathrm BT.
\]
取
\[
m=40\times1.66\times10^{-27}\,\mathrm{kg}=6.64\times10^{-26}\,\mathrm{kg},
\qquad T=(350+273)\,\mathrm K=623\,\mathrm K,
\]
得到
\[
\boxed{v_0=360\,\mathrm{m/s}.}
\]

\paragraph{（2）光子散射产生的减速。}
单个光子动量为 $p_{\mathrm{single}}=h/\lambda$。在时间 $\Delta t$ 内，原子共振吸收 $N=\Gamma\Delta t$ 个反向光子，故
\[
a=\frac{F}{m}=\frac{\Gamma h}{\lambda m}
=7.45\times10^4\,\mathrm{m/s^2}.
\]
恒减速度减速至静止的距离为
\[
\boxed{s=\frac{v_0^2}{2a}=0.87\,\mathrm m.}
\]

\paragraph{（3）无磁场时的多普勒失谐。}
相向运动时，原子参考系中的激光频率为
\[
f'=f\sqrt{\frac{1+v/c}{1-v/c}}.
\]
在 $v\ll c$ 下，共振条件 $f'=f_0$ 给出
\[
\Delta f\equiv f_0-f=\frac{v}{c}f_0=\frac{v}{\lambda}.
\]
因此，当 $v=v_0$ 时，激光应降低的频率为
\[
\boxed{\Delta f=\frac{v_0}{\lambda}=537\,\mathrm{MHz}.}
\]

\paragraph{（4）塞曼补偿磁场。}
沿减速管，原子速度满足
\[
v(z)=\sqrt{v_0^2-2az}.
\]
要保持共振，磁场导致的跃迁频率变化必须抵消多普勒失谐的变化。由 $B(0)=0$、$f_0(z)=f_0(0)+\beta B(z)$，得
\[
\boxed{
B(z)=-\frac{v_0-\sqrt{v_0^2-2az}}{\beta\lambda}
=-\frac{v_0}{\beta\lambda}
\left[1-\sqrt{1-\frac{2\Gamma h z}{\lambda m v_0^2}}\right].}
\]
代入题给数据，
\[
\boxed{B(z)=0.0537\left(1-\sqrt{1-1.15z}\right)\,\mathrm T,}
\]
其中 $z$ 以米为单位，$0\le z\le0.869\,\mathrm m$，相应地 $0\le B\le0.0537\,\mathrm T$。

\subsection{转写依据}
题干依据 \texttt{CPHO38-Q05\_transcription\_v1.md}；公式依据 \texttt{CPHO38-Q05\_formula\_check\_v1.csv} 与答案第12至14页；题图依据原题图5a重绘。本题第（1）至（4）问评分分别为8分、12分、8分、12分。
